\documentclass[11pt,a4paper]{article}
\usepackage[margin=1in]{geometry}
\usepackage[T1]{fontenc}
\usepackage[utf8]{inputenc}
\usepackage{lmodern}
\usepackage{array}
\usepackage{booktabs}
\usepackage{longtable}
\usepackage{tabularx}
\usepackage{enumitem}
\usepackage{amsmath}
\usepackage{hyperref}
\usepackage{xcolor}
\usepackage{listings}
\usepackage{float}
\usepackage{graphicx}

\hypersetup{
  colorlinks=true,
  linkcolor=blue!60!black,
  urlcolor=blue!60!black
}

\lstdefinestyle{plain}{
  basicstyle=\ttfamily\small,
  breaklines=true,
  frame=single,
  columns=fullflexible
}

\title{\textbf{ScalaStream: Judge-Ready Technical Report}\\
Video Streaming Systems -- WebSystems \& Machine Learning Competition\\
Ingenium, IIT Indore}
\author{ScalaStream Team}
\date{March 15, 2026}

\begin{document}
\maketitle
\tableofcontents
\newpage

\section{Executive Summary}
ScalaStream is a containerized Video-on-Demand platform designed to satisfy the competition requirements end-to-end: upload, asynchronous transcoding, low-latency playback, metadata interactions, and personalized recommendations using machine learning.

The system is implemented as production-style microservices using Node.js and Python with Docker Compose:
\begin{itemize}[leftmargin=1.2em]
  \item API Gateway/Auth, Video Service, Transcode Worker, Recommendation Service, and Web Frontend.
  \item PostgreSQL for transactional metadata, Redis for queue and hot counters, MinIO for object storage, and Nginx as stream gateway.
\end{itemize}

Most importantly for judging, the recommendation feed is not static rule logic. It is model-driven through a hybrid ML pipeline that learns from watch history, likes, comments, and search behavior.

\section{Competition Requirement Interpretation}
From the official problem statement, the mandatory pillars are:
\begin{enumerate}[leftmargin=1.4em]
  \item Functional streaming workflow (upload to playback plus metadata).
  \item Non-functional engineering quality (latency, fault handling, scalability/cost tradeoffs).
  \item Real ML-based recommendation (not rule-based), with clear explanation of why ML is used, what data is used, and how UX improves.
\end{enumerate}

This report maps each requirement to implementation evidence and explains technical decisions and tradeoffs.

\section{System Architecture}
\subsection{Component Overview}
\begin{table}[H]
\centering
\begin{tabularx}{\textwidth}{@{}p{3.0cm}p{2.3cm}X@{}}
\toprule
\textbf{Component} & \textbf{Stack} & \textbf{Responsibility} \\
\midrule
Frontend & Node + Vanilla JS + HLS.js & Watch flow, search, auth UX, upload form, recommendation rails, playback controls \\
API Gateway/Auth & Node.js (Express) & Register/login/JWT validation, identity propagation, routing to services \\
Video Service & Node.js (Express) & Upload API, metadata APIs, search, history, qualified view and retention logic \\
Transcode Worker & Python + FFmpeg & Async event consumption, HLS generation (360p/720p), thumbnail extraction, retries \\
Recommendation Service & FastAPI + NumPy & Scheduled training, ranking endpoints, model persistence \\
PostgreSQL & Relational DB & Users, videos, likes, comments, views, watch sessions, search events \\
Redis & Queue + Cache & Stream events for transcoding, hot aggregate counters \\
MinIO & Object Storage & Raw uploads and processed HLS artifacts \\
Nginx Stream Gateway & Nginx & Low-latency delivery of manifests and segments from MinIO \\
\bottomrule
\end{tabularx}
\end{table}

\subsection{Upload to Playback Data Flow}
\begin{enumerate}[leftmargin=1.4em]
  \item User uploads via \texttt{POST /videos/upload}.
  \item Raw file is saved to MinIO (\texttt{raw-videos} bucket), video status becomes \texttt{UPLOADED}.
  \item Event \texttt{video\_uploaded} is published to Redis stream.
  \item Worker consumes event, validates media, transcodes to HLS variants (\texttt{360p}, \texttt{720p}), writes \texttt{master.m3u8}, extracts thumbnail.
  \item Processed files are stored in MinIO (\texttt{processed-videos}), DB status becomes \texttt{READY}.
  \item Frontend streams through \texttt{/stream/\{videoId\}/master.m3u8} served by Nginx.
\end{enumerate}

\subsection{Metadata and Recommendation Data Flow}
\begin{enumerate}[leftmargin=1.4em]
  \item User interactions (likes, comments, views, watch progress, searches) are persisted.
  \item Video aggregates are updated in DB and cached in Redis.
  \item Recommendation service periodically retrains from interaction and content data.
  \item \texttt{GET /feed/recommended} returns per-user ranked feed with explanation tags.
\end{enumerate}

\section{Functional Requirement Compliance}
\begin{longtable}{@{}p{4.5cm}p{2.2cm}p{7.8cm}@{}}
\toprule
\textbf{Requirement} & \textbf{Status} & \textbf{Evidence} \\
\midrule
\endhead
User authentication and authorization & Met & \texttt{POST /auth/register}, \texttt{POST /auth/login}, \texttt{GET /auth/me}; JWT validation in API gateway. \\
Video upload functionality & Met & \texttt{POST /videos/upload} with multipart handling, validation, owner tracking. \\
Video transcoding and storage & Met & Worker creates HLS renditions and stores processed assets in MinIO. \\
Video streaming playback & Met & Dedicated watch route \texttt{/watch/:videoId}; HLS playback via stream gateway. \\
Metadata management (likes/comments/views) & Met & \texttt{/videos/:id/like}, \texttt{/videos/:id/comment}, \texttt{/videos/:id/view} with persistence and counters. \\
ML-based recommendation feed & Met & \texttt{GET /feed/recommended?userId=...} served by trained hybrid model. \\
Concurrent uploads support & Met & Upload ingest decoupled from transcoding using Redis stream queue. \\
\bottomrule
\end{longtable}

\section{Non-Functional Requirement Coverage}
\begin{longtable}{@{}p{4.4cm}p{2.1cm}p{8.0cm}@{}}
\toprule
\textbf{Requirement} & \textbf{Status} & \textbf{Implementation Notes} \\
\midrule
\endhead
Low-latency video delivery & Met & HLS segment delivery through Nginx, cache-friendly headers, short segment duration. \\
High availability and fault tolerance & Mostly Met (local-competition scale) & Queue decoupling, retry and pending-message reclaim in worker, health endpoints, graceful service degradation in UI. \\
Cost-efficient storage/network utilization & Met & Two-rendition ladder (360p/720p), object storage separation, scheduled training instead of per-request training. \\
Scalable architecture & Met (conceptual horizontal scaling) & Stateless API services, scalable worker group consumers, object storage and queue separation. \\
Graceful failure handling & Met & Transcode failure states, retry/finalize controls, status endpoint with queue visibility. \\
Clear tradeoff justification & Met & Documented in scalability/cost strategy with explicit alternatives and rationale. \\
\bottomrule
\end{longtable}

\section{Machine Learning Recommendation System (Critical Requirement)}
\subsection{Why ML Is Used}
Static or manually curated recommendation logic cannot adapt to changing user interests at session level. ScalaStream uses machine learning so the ranking reflects:
\begin{itemize}[leftmargin=1.2em]
  \item implicit preference strength from viewing depth and completion,
  \item explicit engagement from likes and comments,
  \item intent from search history.
\end{itemize}
This directly addresses the requirement that rule-only recommendation is insufficient.

\subsection{Training Data Used}
The model uses behavior and content data from PostgreSQL:
\begin{itemize}[leftmargin=1.2em]
  \item \texttt{video\_views}: watch seconds and completion rate,
  \item \texttt{video\_likes}: positive affinity,
  \item \texttt{video\_comments}: stronger engagement signal,
  \item \texttt{user\_search\_events}: search intent over time,
  \item \texttt{videos} + \texttt{video\_aggregates}: title, description, recency, engagement context.
\end{itemize}

\subsection{Feature Engineering}
For each user-video interaction, an implicit preference score is formed:
\[
\text{score} = 1.45\cdot\text{watch\_ratio}
+ 1.25\cdot\text{completion}
+ 1.35\cdot\text{like\_flag}
+ 1.10\cdot\text{comment\_flag}
+ 0.65\cdot\text{recency\_weight}
\]

Additional representation features:
\begin{itemize}[leftmargin=1.2em]
  \item Item text vectors from title+description (compact hashed n-gram embedding).
  \item User content profile from watched items plus time-decayed search history.
\end{itemize}

\subsection{Model Architecture}
ScalaStream uses a hybrid recommender:
\begin{enumerate}[leftmargin=1.4em]
  \item \textbf{Collaborative layer}: BPR-style pairwise matrix factorization (learned user/item latent vectors).
  \item \textbf{Content layer}: item/user vectors from textual and search signals.
  \item \textbf{Supervised calibration layer}: logistic calibration model trained on sampled positive/negative user-item pairs to learn blend coefficients.
  \item \textbf{Priors}: popularity and recency for tie-break and cold start robustness.
\end{enumerate}

\subsection{Ranking at Inference Time}
Final ranking is a weighted blend:
\[
\text{final} = w_c\cdot\text{collab} + w_t\cdot\text{content}
+ w_l\cdot\text{live\_profile}
+ w_p\cdot\text{popularity}
+ w_r\cdot\text{recency}
\]
Unlike static hand-tuned coefficients, the blend weights are learned during training through logistic loss minimization on sampled ranking pairs, then used at inference.

\subsection{Why This Is Not Rule-Based}
The personalized endpoint \texttt{/feed/recommended} is model-driven because:
\begin{itemize}[leftmargin=1.2em]
  \item latent vectors are learned via iterative optimization,
  \item ranking uses learned affinity scores in embedding space,
  \item blend coefficients are learned by supervised logistic calibration,
  \item training is periodic and data-dependent,
  \item different users receive different ranked outputs from the same catalog.
\end{itemize}

\subsection{Retrain Button Purpose}
The UI ``Retrain Model'' button triggers \texttt{POST /feed/train}. It forces immediate retraining from latest interactions so new watch/like/search behavior is reflected quickly during live demo without waiting for scheduler interval.

\section{Validation Evidence (Local Demonstration)}
\subsection{Service Health Snapshot}
\begin{lstlisting}[style=plain,caption={Health endpoint checks}]
GET /health (api-gateway): ok
GET /health (video-service): ok
GET /health (recommendation-service): ok, model ready=true
GET /health (stream-gateway): ok
\end{lstlisting}

\subsection{Recommendation Model Health Example}
\begin{lstlisting}[style=plain,caption={Recommendation service health summary}]
model.ready = true
training_summary.users = 7
training_summary.videos = 16
training_summary.interactions = 41
training_summary.search_events = 137
training_summary.mode = hybrid_bpr_content
\end{lstlisting}

\subsection{Evidence of Personalization}
In local verification, two different demo users produced different top-5 ranked lists from \texttt{/feed/recommended}. This demonstrates user-dependent ranking behavior, not static ordering.

\subsection{Engagement Integrity Checks}
\begin{itemize}[leftmargin=1.2em]
  \item Repeated like by the same user on same video returns no additional increment (DB unique constraint).
  \item View counting is sessionized and qualified; early watch event may not count while later sufficient watch progression counts.
  \item Retention metrics are derived from \texttt{video\_watch\_sessions} progression.
\end{itemize}

\subsection{Failure Recovery Demonstration}
\texttt{scripts/failure\_demo.ps1} stops the transcode worker, simulates interrupted processing, restarts worker, and polls status until recovery (\texttt{READY}) or timeout.

\subsection{Concurrent Upload Test}
\texttt{scripts/upload\_concurrency\_test.js} uses k6 with 20 virtual users uploading simultaneously to validate queue-based concurrency handling.

\section{Deployment and Submission Readiness}
\subsection{One-Command Startup}
\begin{lstlisting}[style=plain]
docker compose up -d --build
\end{lstlisting}

\subsection{Persistence Guarantees}
Docker volumes persist state across restarts:
\begin{itemize}[leftmargin=1.2em]
  \item \texttt{postgres\_data}
  \item \texttt{minio\_data}
  \item \texttt{redis\_data}
\end{itemize}
Using \texttt{docker compose down} keeps data. Using \texttt{docker compose down -v} clears all state.

\subsection{Independent Judge Run}
The repository is self-contained for judge environments:
\begin{enumerate}[leftmargin=1.4em]
  \item Clone/unzip repository.
  \item Run Docker Compose.
  \item Open frontend at \texttt{http://localhost:3000}.
  \item Use seeded demo accounts or create new users.
  \item Execute demo script from documentation.
\end{enumerate}

\section{Honest Scope, Risks, and Mitigations}
\subsection{Scope Alignment}
\begin{itemize}[leftmargin=1.2em]
  \item Live streaming is intentionally out of scope (VOD only), aligned with competition scope.
  \item Architecture supports horizontal scalability conceptually; deployment target is local competition demonstration.
\end{itemize}

\subsection{Known Risks}
\begin{itemize}[leftmargin=1.2em]
  \item Single-instance database/cache in local mode is not multi-zone HA.
  \item Performance characteristics vary by local machine resources.
\end{itemize}

\subsection{Mitigations}
\begin{itemize}[leftmargin=1.2em]
  \item Queue decoupling and retries reduce failure impact.
  \item Health endpoints and UI degradation messaging improve diagnosability.
  \item Scale path is documented for worker and stateless service replicas.
\end{itemize}

\section{Final Compliance Verdict}
Based on implementation and local validation:
\begin{itemize}[leftmargin=1.2em]
  \item \textbf{Functional requirements:} Met.
  \item \textbf{AI/ML requirement (must be non-rule-based):} Met.
  \item \textbf{Deliverables:} Met.
  \item \textbf{Optional enhancements:} Met.
  \item \textbf{Non-functional requirements:} Met for competition-scale deployment, with clearly documented tradeoffs and horizontal scaling path.
\end{itemize}

\appendix
\section{API Summary}
\begin{longtable}{@{}p{4.0cm}p{10.5cm}@{}}
\toprule
\textbf{Endpoint} & \textbf{Purpose} \\
\midrule
\endhead
\texttt{POST /auth/register} & Create user account \\
\texttt{POST /auth/login} & Authenticate and receive JWT \\
\texttt{GET /auth/me} & Validate and return current user \\
\texttt{POST /videos/upload} & Upload source video and enqueue transcode \\
\texttt{GET /videos/:id/status} & Poll transcode lifecycle state \\
\texttt{POST /videos/:id/finalize} & Requeue failed or force reprocess \\
\texttt{GET /videos} & List ready videos with metadata \\
\texttt{GET /videos/search} & Search by title/description relevance \\
\texttt{GET /videos/:id} & Video details for watch page \\
\texttt{DELETE /videos/:id} & Owner/admin deletion \\
\texttt{POST /videos/:id/like} & Like video (auth required) \\
\texttt{DELETE /videos/:id/like} & Unlike video (auth required) \\
\texttt{POST /videos/:id/comment} & Add comment (auth required) \\
\texttt{POST /videos/:id/view} & Sessionized watch event and qualified view logic \\
\texttt{GET /videos/history/watch} & Watch history for user \\
\texttt{GET /videos/history/search} & Search history for user \\
\texttt{POST /videos/history/search} & Log search query for user \\
\texttt{GET /feed/recommended} & Personalized ML-ranked feed \\
\texttt{GET /feed/trending} & Trending feed mode \\
\texttt{GET /feed/fresh} & Fresh feed mode \\
\texttt{GET /feed/continue} & Continue watching feed mode \\
\texttt{POST /feed/train} & Manual retraining trigger \\
\bottomrule
\end{longtable}

\section{Quick Demo Script (Under 7 Minutes)}
\begin{enumerate}[leftmargin=1.4em]
  \item Start system with Docker Compose and open frontend.
  \item Register or log in with demo accounts.
  \item Upload a video and show status polling until \texttt{READY}.
  \item Open watch page, verify large playback, quality switch, speed control.
  \item Like and comment as authenticated user and show counter updates.
  \item Show watch history/search history and qualified view behavior.
  \item Show \texttt{For You} feed, trigger retrain, and show ranking update.
\end{enumerate}

\section{How to Compile This Report}
\begin{lstlisting}[style=plain]
cd docs
pdflatex ScalaStream_Judge_Report.tex
pdflatex ScalaStream_Judge_Report.tex
\end{lstlisting}

\end{document}
